\chapter{Schätzer}

\section{Ziel dieses Kapitels}

In den vorherigen Kapiteln haben wir Wahrscheinlichkeiten, Wahrscheinlichkeitsdichten, Erwartungswerte, Varianz und Kovarianz kennengelernt.
Jetzt betrachten wir die umgekehrte Situation:

\begin{center}
\textbf{Wir kennen die Verteilung nicht vollständig, sondern haben nur Daten.}
\end{center}

Aus diesen Daten möchten wir unbekannte Größen einer Verteilung bestimmen.
Solche unbekannten Größen nennt man Parameter (parameters).
Eine Vorschrift, die aus Daten einen Parameter schätzt, nennt man Schätzer (estimator).

\begin{conceptbox}
Ein Schätzer (estimator) ist eine Funktion der Daten, mit der wir eine unbekannte Größe einer Verteilung approximieren.

Typische unbekannte Größen sind:

\begin{itemize}
    \item ein Mittelwert \(\mu\),
    \item eine Varianz \(\sigma^2\),
    \item ein Parameter einer Gleichverteilung,
    \item ein Parameter einer Bernoulli-Verteilung,
    \item die Gewichte eines Regressionsmodells.
\end{itemize}
\end{conceptbox}

Nach diesem Kapitel solltest du folgende Fragen beantworten können:

\begin{itemize}
    \item Was ist ein statistischer Schätzer (statistical estimator)?
    \item Was ist der Unterschied zwischen Parameter, Schätzer und Schätzwert?
    \item Was bedeutet i.i.d. (independent and identically distributed)?
    \item Was ist Bias eines Schätzers?
    \item Was ist die Varianz eines Schätzers?
    \item Was ist der mittlere quadratische Fehler (mean squared error)?
    \item Wie hängen Bias, Varianz und Fehler zusammen?
    \item Warum ist das Stichprobenmittel ein erwartungstreuer Schätzer?
    \item Warum ist die naive Stichprobenvarianz verzerrt?
    \item Warum teilt man bei der unverzerrten Stichprobenvarianz durch \(N-1\)?
    \item Was ist Maximum-Likelihood-Schätzung (maximum likelihood estimation)?
    \item Was ist Maximum-a-posteriori-Schätzung (maximum a-posteriori estimation)?
    \item Was ist ein Bayesian Estimator?
\end{itemize}

\section{Parameter, Schätzer und Schätzwert}

Wir beginnen mit einer wichtigen sprachlichen Unterscheidung.

\begin{definitionbox}
Ein Parameter (parameter) ist eine feste, aber unbekannte Größe eines Modells oder einer Verteilung.

Ein Schätzer (estimator) ist eine Funktion der Daten, die diesen Parameter schätzen soll.

Ein Schätzwert (estimate) ist der konkrete Zahlenwert, den der Schätzer für einen konkreten Datensatz liefert.
\end{definitionbox}

Angenommen, die Zufallsvariable \(X\) hat einen unbekannten Mittelwert

\[
    \mu_X = \mathbb{E}[X].
\]

Wir beobachten Daten

\[
    x_1,x_2,\dots,x_N.
\]

Ein möglicher Schätzer für \(\mu_X\) ist das Stichprobenmittel (sample mean):

\[
    \hat{\mu}_N
    =
    \frac{1}{N}
    \sum_{n=1}^{N}
    X_n.
\]

Wichtig ist:
Der Schätzer \(\hat{\mu}_N\) ist selbst eine Zufallsvariable, weil er von den Zufallsvariablen \(X_1,\dots,X_N\) abhängt.
Wenn wir konkrete Daten \(x_1,\dots,x_N\) einsetzen, erhalten wir einen konkreten Schätzwert:

\[
    \hat{\mu}_N(x_1,\dots,x_N)
    =
    \frac{1}{N}
    \sum_{n=1}^{N}
    x_n.
\]

\begin{notebox}
In der Praxis schreibt man oft einfach \(\hat{\mu}_N\), obwohl eigentlich eine Funktion der Daten gemeint ist.

Strenger wäre:

\[
    \hat{\mu}_N
    =
    \hat{\mu}_N(X_1,\dots,X_N).
\]

Diese Unterscheidung ist wichtig:
Der wahre Parameter \(\mu_X\) ist eine feste Zahl.
Der Schätzer \(\hat{\mu}_N\) ist eine Zufallsvariable.
\end{notebox}

\section{Daten als Zufallsvariablen}

In der Statistik und im maschinellen Lernen nehmen wir häufig an, dass Datenpunkte zufällige Realisierungen einer Verteilung sind.

Wir schreiben:

\[
    X_1,X_2,\dots,X_N.
\]

Das sind Zufallsvariablen.
Eine konkrete Beobachtung schreiben wir als

\[
    x_1,x_2,\dots,x_N.
\]

\begin{definitionbox}
Die Daten \(X_1,\dots,X_N\) heißen unabhängig und identisch verteilt (independent and identically distributed), kurz i.i.d., wenn gilt:

\begin{itemize}
    \item Alle \(X_n\) haben dieselbe Verteilung.
    \item Die Zufallsvariablen \(X_1,\dots,X_N\) sind unabhängig.
\end{itemize}
\end{definitionbox}

Identisch verteilt bedeutet:

\[
    X_n \sim p(x)
    \quad
    \text{für alle}
    \quad
    n \in \{1,\dots,N\}.
\]

Unabhängig bedeutet:

\[
    p(x_1,\dots,x_N)
    =
    \prod_{n=1}^{N}
    p(x_n).
\]

\begin{conceptbox}
Die i.i.d.-Annahme ist eine Standardannahme in vielen Herleitungen.

Sie bedeutet:
Alle Datenpunkte stammen aus derselben Quelle und beeinflussen sich nicht gegenseitig.
\end{conceptbox}

\begin{examplebox}
Wenn wir \(N=100\) zufällig ausgewählte Personen messen und \(X_n\) die Körpergröße der \(n\)-ten Person ist, dann modellieren wir oft:

\[
    X_1,\dots,X_N
    \overset{\mathrm{i.i.d.}}{\sim}
    p(x).
\]

Das bedeutet:
Alle Körpergrößen werden als unabhängige Ziehungen aus derselben Verteilung betrachtet.
\end{examplebox}

\section{Beispiel: Parameter einer Gleichverteilung schätzen}

Ein einfaches Beispiel ist die Gleichverteilung (uniform distribution).

Angenommen, die Daten stammen aus einer Gleichverteilung auf einem unbekannten Intervall:

\[
    X \sim \mathrm{Uniform}(a,b).
\]

Die Dichte ist:

\[
    p(x \mid a,b)
    =
    \begin{cases}
        \frac{1}{b-a}, & x \in [a,b], \\
        0, & \text{sonst}.
    \end{cases}
\]

Die Parameter sind \(a\) und \(b\).
Wir beobachten Daten:

\[
    x_1,x_2,\dots,x_N.
\]

Die Frage lautet:

\begin{center}
\textbf{Wie können wir aus den Daten \(a\) und \(b\) schätzen?}
\end{center}

\section{Heuristischer Schätzer mit Minimum und Maximum}

Eine naheliegende Idee ist:
Das kleinste beobachtete \(x_n\) sollte ungefähr bei \(a\) liegen, und das größte beobachtete \(x_n\) ungefähr bei \(b\).

Wir definieren:

\[
    \hat{a}_{\mathrm{mima}}
    =
    \min_{n \in \{1,\dots,N\}} X_n,
\]

und

\[
    \hat{b}_{\mathrm{mima}}
    =
    \max_{n \in \{1,\dots,N\}} X_n.
\]

Der Name \(\mathrm{mima}\) steht hier für Minimum-Maximum.

\begin{definitionbox}
Der Minimum-Maximum-Schätzer (minimum-maximum estimator) für die Gleichverteilung ist:

\[
    \hat{a}_{\mathrm{mima}}
    =
    \min_n X_n,
    \qquad
    \hat{b}_{\mathrm{mima}}
    =
    \max_n X_n.
\]
\end{definitionbox}

\begin{conceptbox}
Die Idee ist sehr intuitiv:
Wenn die Daten aus dem Intervall \([a,b]\) kommen, dann muss \(a\) links vom kleinsten Datenpunkt und \(b\) rechts vom größten Datenpunkt liegen.

Daher verwendet man den kleinsten und größten Datenpunkt als Schätzung der Intervallgrenzen.
\end{conceptbox}

Dieser Schätzer ist einfach, aber er ist nicht perfekt.
Bei endlichem \(N\) ist der kleinste Datenpunkt typischerweise größer als \(a\), und der größte Datenpunkt typischerweise kleiner als \(b\).
Das Intervall wird also tendenziell zu klein geschätzt.

\section{Methode der Momente}

Eine zweite Idee ist die Methode der Momente (method of moments).
Sie nutzt bekannte Formeln für Erwartungswert und Varianz einer Verteilung.

Für die Gleichverteilung auf \([a,b]\) gilt:

\[
    \mathbb{E}[X]
    =
    \frac{a+b}{2},
\]

und

\[
    \mathrm{Var}(X)
    =
    \frac{(b-a)^2}{12}.
\]

Aus Daten können wir Mittelwert und Varianz schätzen.
Dann lösen wir diese Gleichungen nach \(a\) und \(b\) auf.

\subsection{Empirischer Mittelwert}

Der empirische Mittelwert (empirical mean) ist:

\[
    \hat{\mu}_N
    =
    \frac{1}{N}
    \sum_{n=1}^{N}
    X_n.
\]

\subsection{Empirische Varianz}

Eine einfache empirische Varianz ist:

\[
    \hat{\sigma}^2_N
    =
    \frac{1}{N}
    \sum_{n=1}^{N}
    (X_n-\hat{\mu}_N)^2.
\]

Für die Gleichverteilung gilt:

\[
    \mu
    =
    \frac{a+b}{2},
\]

und

\[
    \sigma^2
    =
    \frac{(b-a)^2}{12}.
\]

Daraus folgt:

\[
    a+b = 2\mu,
\]

und

\[
    b-a = \sqrt{12\sigma^2}.
\]

Wir setzen

\[
    d = b-a.
\]

Dann gilt:

\[
    d = \sqrt{12\sigma^2}.
\]

Außerdem:

\[
    a = \mu - \frac{d}{2},
    \qquad
    b = \mu + \frac{d}{2}.
\]

Einsetzen von \(d=\sqrt{12\sigma^2}\) ergibt:

\[
    a = \mu - \sqrt{3\sigma^2},
\]

und

\[
    b = \mu + \sqrt{3\sigma^2}.
\]

Für die Schätzer verwenden wir \(\hat{\mu}_N\) und \(\hat{\sigma}^2_N\):

\[
    \hat{a}_{\mathrm{MoM}}
    =
    \hat{\mu}_N
    -
    \sqrt{3\hat{\sigma}^2_N},
\]

und

\[
    \hat{b}_{\mathrm{MoM}}
    =
    \hat{\mu}_N
    +
    \sqrt{3\hat{\sigma}^2_N}.
\]

\begin{definitionbox}
Die Methode der Momente (method of moments) schätzt Parameter, indem theoretische Momente einer Verteilung mit empirischen Momenten aus den Daten gleichgesetzt werden.

Für die Gleichverteilung ergibt sich:

\[
    \hat{a}_{\mathrm{MoM}}
    =
    \hat{\mu}_N
    -
    \sqrt{3\hat{\sigma}^2_N},
\]

\[
    \hat{b}_{\mathrm{MoM}}
    =
    \hat{\mu}_N
    +
    \sqrt{3\hat{\sigma}^2_N}.
\]
\end{definitionbox}

\begin{notebox}
Momente (moments) sind Erwartungswerte von Potenzen einer Zufallsvariable.

Das erste Moment hängt mit dem Mittelwert zusammen.
Das zweite zentrale Moment ist die Varianz.
\end{notebox}

\section{Welcher Schätzer ist besser?}

Für die Gleichverteilung haben wir nun mindestens zwei Möglichkeiten:

\[
    (\hat{a}_{\mathrm{mima}},\hat{b}_{\mathrm{mima}})
\]

und

\[
    (\hat{a}_{\mathrm{MoM}},\hat{b}_{\mathrm{MoM}}).
\]

Die Frage ist:

\begin{center}
\textbf{Welcher Schätzer ist besser?}
\end{center}

Um das mathematisch zu beantworten, brauchen wir Kriterien.
Die wichtigsten Kriterien sind:

\begin{itemize}
    \item Bias,
    \item Varianz,
    \item mittlerer quadratischer Fehler (mean squared error),
    \item Konsistenz (consistency).
\end{itemize}

\section{Bias eines Schätzers}

Ein Schätzer kann systematisch zu groß oder zu klein sein.
Diese systematische Abweichung nennt man Bias.

\begin{definitionbox}
Sei \(\hat{\theta}_N\) ein Schätzer für einen wahren Parameter \(\theta\).
Der Bias des Schätzers ist

\[
    \mathrm{Bias}[\hat{\theta}_N]
    =
    \mathbb{E}[\hat{\theta}_N] - \theta.
\]
\end{definitionbox}

Falls

\[
    \mathrm{Bias}[\hat{\theta}_N] = 0,
\]

dann heißt der Schätzer erwartungstreu oder unverzerrt (unbiased).

\begin{definitionbox}
Ein Schätzer \(\hat{\theta}_N\) heißt erwartungstreu (unbiased), wenn gilt:

\[
    \mathbb{E}[\hat{\theta}_N] = \theta.
\]
\end{definitionbox}

\begin{conceptbox}
Bias misst systematischen Fehler.

Ein positiver Bias bedeutet:
Der Schätzer ist im Durchschnitt zu groß.

Ein negativer Bias bedeutet:
Der Schätzer ist im Durchschnitt zu klein.
\end{conceptbox}

\section{Varianz eines Schätzers}

Neben systematischem Fehler gibt es zufällige Schwankung.
Ein Schätzer kann bei verschiedenen Datensätzen stark unterschiedliche Werte liefern.
Diese Schwankung misst die Varianz des Schätzers.

\begin{definitionbox}
Die Varianz eines Schätzers \(\hat{\theta}_N\) ist

\[
    \mathrm{Var}[\hat{\theta}_N]
    =
    \mathbb{E}\left[
        \left(
            \hat{\theta}_N
            -
            \mathbb{E}[\hat{\theta}_N]
        \right)^2
    \right].
\]
\end{definitionbox}

\begin{conceptbox}
Die Varianz eines Schätzers misst, wie stark der Schätzer von Datensatz zu Datensatz schwankt.

Kleine Varianz bedeutet:
Der Schätzer ist stabil.

Große Varianz bedeutet:
Der Schätzer reagiert stark auf die zufällige Auswahl der Daten.
\end{conceptbox}

\section{Mittlerer quadratischer Fehler}

Bias und Varianz beschreiben unterschiedliche Fehlerquellen.
Der mittlere quadratische Fehler (mean squared error, MSE) kombiniert beide.

\begin{definitionbox}
Der mittlere quadratische Fehler (mean squared error) eines Schätzers \(\hat{\theta}_N\) für einen Parameter \(\theta\) ist

\[
    \mathrm{MSE}[\hat{\theta}_N]
    =
    \mathbb{E}\left[
        \left(
            \hat{\theta}_N - \theta
        \right)^2
    \right].
\]
\end{definitionbox}

\subsection{Bias-Varianz-Zerlegung}

Der MSE kann zerlegt werden in Bias und Varianz:

\[
    \mathrm{MSE}[\hat{\theta}_N]
    =
    \mathrm{Bias}[\hat{\theta}_N]^2
    +
    \mathrm{Var}[\hat{\theta}_N].
\]

\subsection{Herleitung}

Wir starten mit der Definition:

\[
    \mathrm{MSE}[\hat{\theta}_N]
    =
    \mathbb{E}\left[
        \left(
            \hat{\theta}_N-\theta
        \right)^2
    \right].
\]

Wir addieren und subtrahieren \(\mathbb{E}[\hat{\theta}_N]\):

\[
    \hat{\theta}_N-\theta
    =
    \hat{\theta}_N
    -
    \mathbb{E}[\hat{\theta}_N]
    +
    \mathbb{E}[\hat{\theta}_N]
    -
    \theta.
\]

Setze

\[
    A
    =
    \hat{\theta}_N
    -
    \mathbb{E}[\hat{\theta}_N],
\]

und

\[
    B
    =
    \mathbb{E}[\hat{\theta}_N]
    -
    \theta.
\]

Dann gilt:

\[
    \hat{\theta}_N-\theta
    =
    A+B.
\]

Also:

\[
    \mathrm{MSE}[\hat{\theta}_N]
    =
    \mathbb{E}[(A+B)^2].
\]

Ausmultiplizieren ergibt:

\[
    \mathbb{E}[(A+B)^2]
    =
    \mathbb{E}[A^2]
    +
    2\mathbb{E}[AB]
    +
    \mathbb{E}[B^2].
\]

Da \(B\) konstant ist, gilt:

\[
    \mathbb{E}[AB]
    =
    B\mathbb{E}[A].
\]

Außerdem:

\[
    \mathbb{E}[A]
    =
    \mathbb{E}\left[
        \hat{\theta}_N
        -
        \mathbb{E}[\hat{\theta}_N]
    \right]
    =
    \mathbb{E}[\hat{\theta}_N]
    -
    \mathbb{E}[\hat{\theta}_N]
    =
    0.
\]

Daher verschwindet der Kreuzterm:

\[
    2\mathbb{E}[AB] = 0.
\]

Weiter gilt:

\[
    \mathbb{E}[A^2]
    =
    \mathrm{Var}[\hat{\theta}_N],
\]

und

\[
    \mathbb{E}[B^2]
    =
    B^2
    =
    \left(
        \mathbb{E}[\hat{\theta}_N]
        -
        \theta
    \right)^2
    =
    \mathrm{Bias}[\hat{\theta}_N]^2.
\]

Damit:

\[
    \mathrm{MSE}[\hat{\theta}_N]
    =
    \mathrm{Var}[\hat{\theta}_N]
    +
    \mathrm{Bias}[\hat{\theta}_N]^2.
\]

\begin{conceptbox}
Die Bias-Varianz-Zerlegung lautet:

\[
    \mathrm{MSE}
    =
    \text{Bias}^2
    +
    \text{Varianz}.
\]

Ein guter Schätzer sollte idealerweise kleinen Bias und kleine Varianz haben.
Oft stehen diese beiden Ziele aber in Konkurrenz.
\end{conceptbox}

\section{Bias-Varianz-Trade-off}

Ein Schätzer mit sehr kleinem Bias kann hohe Varianz haben.
Ein Schätzer mit kleiner Varianz kann dafür verzerrt sein.
Dieses Spannungsverhältnis nennt man Bias-Varianz-Trade-off (bias-variance trade-off).

\begin{conceptbox}
Beim Bias-Varianz-Trade-off versucht man, den gesamten erwarteten Fehler zu minimieren, nicht nur Bias oder Varianz einzeln.

Ein leicht verzerrter Schätzer kann besser sein als ein unverzerrter Schätzer, wenn seine Varianz deutlich kleiner ist.
\end{conceptbox}

\begin{examtipbox}
Wichtig:
Unverzerrt bedeutet nicht automatisch gut.

Ein Schätzer kann Bias \(0\) haben, aber sehr stark schwanken.
Dann kann sein MSE trotzdem groß sein.
\end{examtipbox}

\section{Konsistenz}

Ein Schätzer sollte besser werden, wenn mehr Daten verfügbar sind.
Diese Eigenschaft nennt man Konsistenz (consistency).

\begin{definitionbox}
Ein Schätzer \(\hat{\theta}_N\) heißt konsistent (consistent), wenn er für \(N \to \infty\) gegen den wahren Parameter \(\theta\) konvergiert.

Eine einfache ausreichende Intuition ist:

\[
    \lim_{N \to \infty}
    \mathrm{Bias}[\hat{\theta}_N]
    =
    0,
\]

und

\[
    \lim_{N \to \infty}
    \mathrm{Var}[\hat{\theta}_N]
    =
    0.
\]
\end{definitionbox}

\begin{conceptbox}
Konsistenz bedeutet:
Mit unendlich vielen Daten verschwindet der Schätzfehler.

Für maschinelles Lernen ist das wichtig, weil wir erwarten, dass mehr Daten zu besseren Schätzungen führen.
\end{conceptbox}

\section{Das Stichprobenmittel}

Das Stichprobenmittel (sample mean) ist einer der wichtigsten Schätzer.

Seien

\[
    X_1,\dots,X_N
\]

i.i.d. Zufallsvariablen mit

\[
    \mathbb{E}[X_n] = \mu_X
\]

und

\[
    \mathrm{Var}(X_n) = \sigma_X^2.
\]

Der Schätzer für den Mittelwert ist:

\[
    \hat{\mu}_N
    =
    \frac{1}{N}
    \sum_{n=1}^{N}
    X_n.
\]

\begin{definitionbox}
Das Stichprobenmittel (sample mean, empirical mean) ist

\[
    \hat{\mu}_N
    =
    \frac{1}{N}
    \sum_{n=1}^{N}
    X_n.
\]

Es schätzt den wahren Mittelwert

\[
    \mu_X
    =
    \mathbb{E}[X].
\]
\end{definitionbox}

\section{Bias des Stichprobenmittels}

Wir berechnen:

\[
    \mathrm{Bias}[\hat{\mu}_N]
    =
    \mathbb{E}[\hat{\mu}_N]
    -
    \mu_X.
\]

Zuerst:

\[
    \mathbb{E}[\hat{\mu}_N]
    =
    \mathbb{E}\left[
        \frac{1}{N}
        \sum_{n=1}^{N}
        X_n
    \right].
\]

Wegen Linearität des Erwartungswerts:

\[
    \mathbb{E}[\hat{\mu}_N]
    =
    \frac{1}{N}
    \sum_{n=1}^{N}
    \mathbb{E}[X_n].
\]

Da alle \(X_n\) identisch verteilt sind:

\[
    \mathbb{E}[X_n] = \mu_X.
\]

Also:

\[
    \mathbb{E}[\hat{\mu}_N]
    =
    \frac{1}{N}
    \sum_{n=1}^{N}
    \mu_X
    =
    \frac{1}{N}
    N\mu_X
    =
    \mu_X.
\]

Damit:

\[
    \mathrm{Bias}[\hat{\mu}_N]
    =
    \mu_X-\mu_X
    =
    0.
\]

\begin{conceptbox}
Das Stichprobenmittel ist ein erwartungstreuer Schätzer (unbiased estimator) für den wahren Mittelwert:

\[
    \mathbb{E}[\hat{\mu}_N]
    =
    \mu_X.
\]
\end{conceptbox}

\section{Varianz des Stichprobenmittels}

Nun berechnen wir die Varianz von \(\hat{\mu}_N\):

\[
    \mathrm{Var}[\hat{\mu}_N]
    =
    \mathrm{Var}\left[
        \frac{1}{N}
        \sum_{n=1}^{N}
        X_n
    \right].
\]

Konstanten ziehen quadratisch aus der Varianz:

\[
    \mathrm{Var}[\hat{\mu}_N]
    =
    \frac{1}{N^2}
    \mathrm{Var}\left[
        \sum_{n=1}^{N}
        X_n
    \right].
\]

Da die \(X_n\) unabhängig sind, addieren sich die Varianzen:

\[
    \mathrm{Var}\left[
        \sum_{n=1}^{N}
        X_n
    \right]
    =
    \sum_{n=1}^{N}
    \mathrm{Var}[X_n].
\]

Da alle \(X_n\) identisch verteilt sind:

\[
    \mathrm{Var}[X_n] = \sigma_X^2.
\]

Also:

\[
    \sum_{n=1}^{N}
    \mathrm{Var}[X_n]
    =
    N\sigma_X^2.
\]

Damit:

\[
    \mathrm{Var}[\hat{\mu}_N]
    =
    \frac{1}{N^2}
    N\sigma_X^2
    =
    \frac{\sigma_X^2}{N}.
\]

\begin{definitionbox}
Für das Stichprobenmittel gilt:

\[
    \mathrm{Bias}[\hat{\mu}_N] = 0,
\]

und

\[
    \mathrm{Var}[\hat{\mu}_N]
    =
    \frac{\sigma_X^2}{N}.
\]
\end{definitionbox}

\begin{conceptbox}
Je größer \(N\) ist, desto kleiner ist die Varianz des Stichprobenmittels.

\[
    \mathrm{Var}[\hat{\mu}_N]
    =
    \frac{\sigma_X^2}{N}
    \longrightarrow
    0
    \quad
    \text{für}
    \quad
    N \to \infty.
\]

Das Stichprobenmittel wird also mit mehr Daten stabiler.
\end{conceptbox}

Da der Bias \(0\) ist, ist der MSE des Stichprobenmittels gleich seiner Varianz:

\[
    \mathrm{MSE}[\hat{\mu}_N]
    =
    \mathrm{Var}[\hat{\mu}_N]
    =
    \frac{\sigma_X^2}{N}.
\]

\section{Die naive Stichprobenvarianz}

Wir möchten nun die wahre Varianz

\[
    \sigma_X^2
    =
    \mathrm{Var}(X)
\]

schätzen.

Eine naheliegende Idee ist:

\[
    \hat{\sigma}^2_N
    =
    \frac{1}{N}
    \sum_{n=1}^{N}
    (X_n-\hat{\mu}_N)^2.
\]

\begin{definitionbox}
Die naive Stichprobenvarianz ist

\[
    \hat{\sigma}^2_N
    =
    \frac{1}{N}
    \sum_{n=1}^{N}
    (X_n-\hat{\mu}_N)^2.
\]

Sie verwendet das Stichprobenmittel \(\hat{\mu}_N\) anstelle des unbekannten wahren Mittelwerts \(\mu_X\).
\end{definitionbox}

Diese Formel wirkt natürlich, ist aber verzerrt.
Sie unterschätzt die wahre Varianz im Mittel.

\section{Bias der naiven Stichprobenvarianz}

Für die naive Stichprobenvarianz gilt:

\[
    \mathbb{E}[\hat{\sigma}^2_N]
    =
    \frac{N-1}{N}
    \sigma_X^2.
\]

Damit ist der Bias:

\[
    \mathrm{Bias}[\hat{\sigma}^2_N]
    =
    \mathbb{E}[\hat{\sigma}^2_N]
    -
    \sigma_X^2.
\]

Einsetzen ergibt:

\[
    \mathrm{Bias}[\hat{\sigma}^2_N]
    =
    \frac{N-1}{N}
    \sigma_X^2
    -
    \sigma_X^2.
\]

Wir schreiben

\[
    \sigma_X^2
    =
    \frac{N}{N}
    \sigma_X^2.
\]

Dann:

\[
    \mathrm{Bias}[\hat{\sigma}^2_N]
    =
    \frac{N-1}{N}
    \sigma_X^2
    -
    \frac{N}{N}
    \sigma_X^2
    =
    -\frac{1}{N}
    \sigma_X^2.
\]

\begin{conceptbox}
Die naive Stichprobenvarianz ist verzerrt:

\[
    \mathrm{Bias}[\hat{\sigma}^2_N]
    =
    -\frac{1}{N}
    \sigma_X^2.
\]

Sie unterschätzt die wahre Varianz im Durchschnitt.
\end{conceptbox}

\subsection{Intuition für die Verzerrung}

Der Grund ist:
Wir kennen den wahren Mittelwert \(\mu_X\) nicht.
Stattdessen verwenden wir das aus denselben Daten berechnete Stichprobenmittel \(\hat{\mu}_N\).

Das Stichprobenmittel liegt näher an den Daten als der wahre Mittelwert.
Dadurch werden die Abstände

\[
    X_n-\hat{\mu}_N
\]

im Durchschnitt zu klein.
Also wird auch die Varianz zu klein geschätzt.

\section{Unverzerrte Stichprobenvarianz}

Um den Bias zu korrigieren, multiplizieren wir die naive Stichprobenvarianz mit

\[
    \frac{N}{N-1}.
\]

Wir definieren:

\[
    \hat{s}^2_N
    =
    \frac{N}{N-1}
    \hat{\sigma}^2_N.
\]

Einsetzen der naiven Stichprobenvarianz ergibt:

\[
    \hat{s}^2_N
    =
    \frac{N}{N-1}
    \frac{1}{N}
    \sum_{n=1}^{N}
    (X_n-\hat{\mu}_N)^2.
\]

Damit:

\[
    \hat{s}^2_N
    =
    \frac{1}{N-1}
    \sum_{n=1}^{N}
    (X_n-\hat{\mu}_N)^2.
\]

\begin{definitionbox}
Die unverzerrte Stichprobenvarianz ist

\[
    \hat{s}^2_N
    =
    \frac{1}{N-1}
    \sum_{n=1}^{N}
    (X_n-\hat{\mu}_N)^2.
\]

Sie ist ein erwartungstreuer Schätzer für \(\sigma_X^2\).
\end{definitionbox}

\subsection{Nachweis der Erwartungstreue}

Wir wissen:

\[
    \mathbb{E}[\hat{\sigma}^2_N]
    =
    \frac{N-1}{N}
    \sigma_X^2.
\]

Da

\[
    \hat{s}^2_N
    =
    \frac{N}{N-1}
    \hat{\sigma}^2_N
\]

gilt:

\[
    \mathbb{E}[\hat{s}^2_N]
    =
    \mathbb{E}\left[
        \frac{N}{N-1}
        \hat{\sigma}^2_N
    \right].
\]

Der Faktor ist konstant:

\[
    \mathbb{E}[\hat{s}^2_N]
    =
    \frac{N}{N-1}
    \mathbb{E}[\hat{\sigma}^2_N].
\]

Einsetzen:

\[
    \mathbb{E}[\hat{s}^2_N]
    =
    \frac{N}{N-1}
    \frac{N-1}{N}
    \sigma_X^2.
\]

Damit:

\[
    \mathbb{E}[\hat{s}^2_N]
    =
    \sigma_X^2.
\]

Also:

\[
    \mathrm{Bias}[\hat{s}^2_N] = 0.
\]

\begin{examtipbox}
Der Grund für \(N-1\) im Nenner ist die Bias-Korrektur.

\[
    \frac{1}{N}
    \sum_{n=1}^{N}
    (X_n-\hat{\mu}_N)^2
\]

ist verzerrt.

\[
    \frac{1}{N-1}
    \sum_{n=1}^{N}
    (X_n-\hat{\mu}_N)^2
\]

ist unverzerrt.
\end{examtipbox}

\section{Vergleich der Varianzschätzer}

Wir haben zwei Schätzer für \(\sigma_X^2\):

\[
    \hat{\sigma}^2_N
    =
    \frac{1}{N}
    \sum_{n=1}^{N}
    (X_n-\hat{\mu}_N)^2,
\]

und

\[
    \hat{s}^2_N
    =
    \frac{1}{N-1}
    \sum_{n=1}^{N}
    (X_n-\hat{\mu}_N)^2.
\]

Zwischen ihnen gilt:

\[
    \hat{s}^2_N
    =
    \frac{N}{N-1}
    \hat{\sigma}^2_N.
\]

Daher gilt für die Varianzen:

\[
    \mathrm{Var}[\hat{s}^2_N]
    =
    \mathrm{Var}\left[
        \frac{N}{N-1}
        \hat{\sigma}^2_N
    \right].
\]

Konstanten gehen quadratisch in die Varianz ein:

\[
    \mathrm{Var}[\hat{s}^2_N]
    =
    \left(
        \frac{N}{N-1}
    \right)^2
    \mathrm{Var}[\hat{\sigma}^2_N].
\]

Da

\[
    \frac{N}{N-1} > 1
\]

für \(N>1\), gilt:

\[
    \mathrm{Var}[\hat{s}^2_N]
    >
    \mathrm{Var}[\hat{\sigma}^2_N].
\]

\begin{conceptbox}
Die unverzerrte Stichprobenvarianz hat Bias \(0\), aber etwas größere Varianz als die naive Stichprobenvarianz.

Das ist ein Beispiel für den Bias-Varianz-Trade-off.
\end{conceptbox}

\section{Übersicht: Mittelwert- und Varianzschätzer}

\[
\begin{array}{c|c|c|c}
\text{Zielgröße}
&
\text{Schätzer}
&
\text{Bias}
&
\text{Varianz}
\\
\hline
\mu_X
&
\hat{\mu}_N
=
\frac{1}{N}\sum_{n=1}^{N}X_n
&
0
&
\frac{1}{N}\sigma_X^2
\\
\hline
\sigma_X^2
&
\hat{\sigma}^2_N
=
\frac{1}{N}
\sum_{n=1}^{N}
(X_n-\hat{\mu}_N)^2
&
-\frac{1}{N}\sigma_X^2
&
\mathrm{Var}[\hat{\sigma}^2_N]
\\
\hline
\sigma_X^2
&
\hat{s}^2_N
=
\frac{1}{N-1}
\sum_{n=1}^{N}
(X_n-\hat{\mu}_N)^2
&
0
&
\left(\frac{N}{N-1}\right)^2
\mathrm{Var}[\hat{\sigma}^2_N]
\end{array}
\]

\begin{notebox}
Die Tabelle zeigt:
Der unverzerrte Varianzschätzer reduziert den Bias auf \(0\), zahlt dafür aber mit etwas höherer Varianz.
Für große \(N\) wird der Unterschied zwischen \(N\) und \(N-1\) klein.
\end{notebox}

\section{Maximum-Likelihood-Schätzung}

Die Maximum-Likelihood-Schätzung (maximum likelihood estimation, MLE) ist eine der wichtigsten Methoden im maschinellen Lernen.

Die Idee ist:

\begin{center}
\textbf{Wähle die Parameter so, dass die beobachteten Daten unter dem Modell möglichst wahrscheinlich sind.}
\end{center}

\begin{definitionbox}
Sei \(p(\vec{X} \mid \theta)\) die Wahrscheinlichkeit oder Dichte der Daten unter dem Parameter \(\theta\).

Die Likelihood ist

\[
    L(\theta;\vec{X})
    =
    p(\vec{X} \mid \theta).
\]

Der Maximum-Likelihood-Schätzer ist

\[
    \hat{\theta}_{\mathrm{ML}}
    =
    \arg\max_{\theta}
    p(\vec{X} \mid \theta).
\]
\end{definitionbox}

Hier bezeichnet

\[
    \vec{X}
    =
    (X_1,\dots,X_N)^\top
\]

den gesamten Datensatz als Zufallsvektor.

Für konkrete Daten

\[
    \vec{x}
    =
    (x_1,\dots,x_N)^\top
\]

betrachtet man

\[
    L(\theta;\vec{x})
    =
    p(\vec{x} \mid \theta)
\]

als Funktion von \(\theta\).

\begin{notebox}
Die Likelihood \(L(\theta;\vec{x})\) ist keine Wahrscheinlichkeitsverteilung über \(\theta\).

Sie ist eine Funktion, die sagt, wie plausibel die beobachteten Daten \(\vec{x}\) für verschiedene Parameterwerte \(\theta\) sind.
\end{notebox}

\section{Likelihood bei i.i.d. Daten}

Wenn die Daten i.i.d. sind, gilt:

\[
    p(\vec{x} \mid \theta)
    =
    p(x_1,\dots,x_N \mid \theta)
    =
    \prod_{n=1}^{N}
    p(x_n \mid \theta).
\]

Damit ist der Maximum-Likelihood-Schätzer:

\[
    \hat{\theta}_{\mathrm{ML}}
    =
    \arg\max_{\theta}
    \prod_{n=1}^{N}
    p(x_n \mid \theta).
\]

\begin{conceptbox}
Bei i.i.d. Daten faktorisiert die Likelihood als Produkt der Einzelwahrscheinlichkeiten oder Einzeldichten:

\[
    L(\theta;\vec{x})
    =
    \prod_{n=1}^{N}
    p(x_n \mid \theta).
\]
\end{conceptbox}

\section{Log-Likelihood}

Produkte vieler Wahrscheinlichkeiten sind mathematisch unpraktisch.
Deshalb nimmt man meistens den Logarithmus.
Da der Logarithmus streng monoton wächst, ändert sich die Stelle des Maximums nicht.

\[
    \hat{\theta}_{\mathrm{ML}}
    =
    \arg\max_{\theta}
    L(\theta;\vec{x})
\]

ist äquivalent zu

\[
    \hat{\theta}_{\mathrm{ML}}
    =
    \arg\max_{\theta}
    \log L(\theta;\vec{x}).
\]

Für i.i.d. Daten gilt:

\[
    \log L(\theta;\vec{x})
    =
    \log
    \prod_{n=1}^{N}
    p(x_n \mid \theta).
\]

Der Logarithmus macht aus dem Produkt eine Summe:

\[
    \log L(\theta;\vec{x})
    =
    \sum_{n=1}^{N}
    \log p(x_n \mid \theta).
\]

\begin{definitionbox}
Die Log-Likelihood ist

\[
    \ell(\theta)
    =
    \log L(\theta;\vec{x})
    =
    \sum_{n=1}^{N}
    \log p(x_n \mid \theta).
\]
\end{definitionbox}

\begin{examtipbox}
Standardtechnik für Maximum-Likelihood-Herleitungen:

\begin{enumerate}
    \item Likelihood \(L(\theta;\vec{x})\) aufschreiben.
    \item Log-Likelihood \(\ell(\theta)\) bilden.
    \item Nach \(\theta\) ableiten.
    \item Ableitung gleich \(0\) setzen.
    \item Nach \(\theta\) lösen.
    \item Prüfen, ob die Lösung zulässig ist und wirklich ein Maximum liefert.
\end{enumerate}
\end{examtipbox}

\section{Beispiel: MLE für Bernoulli-Verteilung}

Die Bernoulli-Verteilung (Bernoulli distribution) modelliert binäre Daten:

\[
    X \in \{0,1\}.
\]

Der Parameter sei

\[
    \pi \in [0,1].
\]

Die Verteilung ist:

\[
    p(X=x \mid \pi)
    =
    \pi^x(1-\pi)^{1-x}.
\]

Wir beobachten i.i.d. Daten:

\[
    x_1,\dots,x_N
    \in
    \{0,1\}.
\]

Die Likelihood ist:

\[
    L(\pi;\vec{x})
    =
    \prod_{n=1}^{N}
    \pi^{x_n}
    (1-\pi)^{1-x_n}.
\]

Die Log-Likelihood ist:

\[
    \ell(\pi)
    =
    \sum_{n=1}^{N}
    \left[
        x_n \log \pi
        +
        (1-x_n)\log(1-\pi)
    \right].
\]

Ableiten nach \(\pi\):

\[
    \frac{\partial \ell}{\partial \pi}
    =
    \sum_{n=1}^{N}
    \left[
        \frac{x_n}{\pi}
        -
        \frac{1-x_n}{1-\pi}
    \right].
\]

Setze die Ableitung gleich \(0\):

\[
    \sum_{n=1}^{N}
    \left[
        \frac{x_n}{\pi}
        -
        \frac{1-x_n}{1-\pi}
    \right]
    =
    0.
\]

Das ist:

\[
    \frac{1}{\pi}
    \sum_{n=1}^{N}
    x_n
    -
    \frac{1}{1-\pi}
    \sum_{n=1}^{N}
    (1-x_n)
    =
    0.
\]

Setze

\[
    S
    =
    \sum_{n=1}^{N}
    x_n.
\]

Dann ist

\[
    \sum_{n=1}^{N}
    (1-x_n)
    =
    N-S.
\]

Also:

\[
    \frac{S}{\pi}
    -
    \frac{N-S}{1-\pi}
    =
    0.
\]

Damit:

\[
    \frac{S}{\pi}
    =
    \frac{N-S}{1-\pi}.
\]

Kreuzmultiplizieren:

\[
    S(1-\pi)
    =
    \pi(N-S).
\]

Ausmultiplizieren:

\[
    S - S\pi
    =
    \pi N - \pi S.
\]

Die Terme \(-S\pi\) und \(-\pi S\) heben sich auf:

\[
    S
    =
    \pi N.
\]

Also:

\[
    \hat{\pi}_{\mathrm{ML}}
    =
    \frac{S}{N}
    =
    \frac{1}{N}
    \sum_{n=1}^{N}
    x_n.
\]

\begin{conceptbox}
Der Maximum-Likelihood-Schätzer für den Parameter einer Bernoulli-Verteilung ist der Anteil der Einsen:

\[
    \hat{\pi}_{\mathrm{ML}}
    =
    \frac{1}{N}
    \sum_{n=1}^{N}
    x_n.
\]
\end{conceptbox}

\section{MLE für die Normalverteilung}

Seien

\[
    X_1,\dots,X_N
    \overset{\mathrm{i.i.d.}}{\sim}
    \mathcal{N}(\mu,\sigma^2).
\]

Die Parameter sind \(\mu\) und \(\sigma^2\).
Für die Normalverteilung ergibt Maximum-Likelihood:

\[
    \hat{\mu}_{\mathrm{ML}}
    =
    \frac{1}{N}
    \sum_{n=1}^{N}
    x_n,
\]

und

\[
    \hat{\sigma}^2_{\mathrm{ML}}
    =
    \frac{1}{N}
    \sum_{n=1}^{N}
    (x_n-\hat{\mu}_{\mathrm{ML}})^2.
\]

\begin{notebox}
Der ML-Schätzer für die Varianz der Normalverteilung teilt durch \(N\), nicht durch \(N-1\).

Er ist daher identisch mit der naiven Stichprobenvarianz und ist bei endlichem \(N\) verzerrt.
\end{notebox}

\begin{conceptbox}
Maximum-Likelihood-Schätzer sind oft sehr allgemein herleitbar.
Sie sind deshalb für maschinelles Lernen besonders wichtig.

Aber:
MLE ist nicht automatisch unverzerrt.
\end{conceptbox}

\section{MLE für die Gleichverteilung}

Für die Gleichverteilung

\[
    X \sim \mathrm{Uniform}(a,b)
\]

ist die Dichte:

\[
    p(x \mid a,b)
    =
    \begin{cases}
        \frac{1}{b-a}, & x \in [a,b], \\
        0, & \text{sonst}.
    \end{cases}
\]

Für i.i.d. Daten ergibt sich die Likelihood:

\[
    L(a,b;\vec{x})
    =
    \prod_{n=1}^{N}
    p(x_n \mid a,b).
\]

Damit gilt:

\[
    L(a,b;\vec{x})
    =
    \begin{cases}
        \left(\frac{1}{b-a}\right)^N,
        & \text{wenn } a \leq x_n \leq b \text{ für alle } n, \\
        0,
        & \text{sonst}.
    \end{cases}
\]

Damit die Likelihood nicht \(0\) ist, muss das Intervall \([a,b]\) alle Datenpunkte enthalten:

\[
    a \leq \min_n x_n,
    \qquad
    b \geq \max_n x_n.
\]

Wenn diese Bedingung erfüllt ist, ist

\[
    L(a,b;\vec{x})
    =
    \left(\frac{1}{b-a}\right)^N.
\]

Diese Likelihood wird größer, wenn \(b-a\) kleiner wird.
Das kleinste Intervall, das alle Daten enthält, ist:

\[
    \hat{a}_{\mathrm{ML}}
    =
    \min_n x_n,
\]

und

\[
    \hat{b}_{\mathrm{ML}}
    =
    \max_n x_n.
\]

\begin{definitionbox}
Für die Gleichverteilung auf \([a,b]\) ist der Maximum-Likelihood-Schätzer:

\[
    \hat{a}_{\mathrm{ML}}
    =
    \min_n x_n,
    \qquad
    \hat{b}_{\mathrm{ML}}
    =
    \max_n x_n.
\]
\end{definitionbox}

\begin{notebox}
Für die Gleichverteilung fällt der Maximum-Likelihood-Schätzer mit dem Minimum-Maximum-Schätzer zusammen.

Das ist ein gutes Beispiel dafür, dass MLE nicht immer alle gewünschten Eigenschaften hat:
Bei endlichem \(N\) schätzt dieser Schätzer das Intervall tendenziell zu klein.
\end{notebox}

\section{Maximum-a-posteriori-Schätzung}

Maximum-Likelihood verwendet nur die Daten.
Manchmal haben wir aber Vorwissen über plausible Parameter.
Dieses Vorwissen modellieren wir mit einem Prior.

\begin{definitionbox}
Bei der Maximum-a-posteriori-Schätzung (maximum a-posteriori estimation, MAP) maximiert man die Posterior-Verteilung der Parameter:

\[
    \hat{\theta}_{\mathrm{MAP}}
    =
    \arg\max_{\theta}
    p(\theta \mid \vec{x}).
\]
\end{definitionbox}

Mit Bayes' Regel gilt:

\[
    p(\theta \mid \vec{x})
    =
    \frac{
        p(\vec{x} \mid \theta)p(\theta)
    }{
        p(\vec{x})
    }.
\]

Dabei ist:

\begin{itemize}
    \item \(p(\vec{x} \mid \theta)\) die Likelihood,
    \item \(p(\theta)\) der Prior,
    \item \(p(\theta \mid \vec{x})\) der Posterior,
    \item \(p(\vec{x})\) die Evidence.
\end{itemize}

Da \(p(\vec{x})\) nicht von \(\theta\) abhängt, ist es für die Maximierung konstant.
Daher:

\[
    \hat{\theta}_{\mathrm{MAP}}
    =
    \arg\max_{\theta}
    p(\vec{x} \mid \theta)p(\theta).
\]

Oder proportional geschrieben:

\[
    p(\theta \mid \vec{x})
    \propto
    p(\vec{x} \mid \theta)p(\theta).
\]

\begin{conceptbox}
MAP kombiniert Daten und Vorwissen:

\[
    \text{Posterior}
    \propto
    \text{Likelihood}
    \cdot
    \text{Prior}.
\]

MLE maximiert nur die Likelihood.
MAP maximiert Likelihood mal Prior.
\end{conceptbox}

\section{Log-Posterior bei MAP}

Wie bei MLE verwendet man oft den Logarithmus:

\[
    \hat{\theta}_{\mathrm{MAP}}
    =
    \arg\max_{\theta}
    \left[
        \log p(\vec{x} \mid \theta)
        +
        \log p(\theta)
    \right].
\]

Für i.i.d. Daten ist:

\[
    \log p(\vec{x} \mid \theta)
    =
    \sum_{n=1}^{N}
    \log p(x_n \mid \theta).
\]

Damit:

\[
    \hat{\theta}_{\mathrm{MAP}}
    =
    \arg\max_{\theta}
    \left[
        \sum_{n=1}^{N}
        \log p(x_n \mid \theta)
        +
        \log p(\theta)
    \right].
\]

\begin{examtipbox}
Standardtechnik für MAP-Herleitungen:

\begin{enumerate}
    \item Likelihood \(p(\vec{x}\mid\theta)\) aufschreiben.
    \item Prior \(p(\theta)\) aufschreiben.
    \item Produkt \(p(\vec{x}\mid\theta)p(\theta)\) bilden.
    \item Logarithmus nehmen.
    \item Nach \(\theta\) ableiten.
    \item Ableitung gleich \(0\) setzen.
    \item Zulässige optimale Lösung wählen.
\end{enumerate}
\end{examtipbox}

\section{MLE als Spezialfall von MAP}

Wenn der Prior konstant ist, dann ist MAP äquivalent zu MLE.

Angenommen:

\[
    p(\theta) = \text{konstant}.
\]

Dann gilt:

\[
    \hat{\theta}_{\mathrm{MAP}}
    =
    \arg\max_{\theta}
    p(\vec{x}\mid\theta)p(\theta).
\]

Da \(p(\theta)\) konstant ist, ändert es die Stelle des Maximums nicht:

\[
    \hat{\theta}_{\mathrm{MAP}}
    =
    \arg\max_{\theta}
    p(\vec{x}\mid\theta)
    =
    \hat{\theta}_{\mathrm{ML}}.
\]

\begin{conceptbox}
MLE ist MAP mit einem konstanten oder uninformativem Prior.

MAP erweitert MLE, indem Vorwissen über plausible Parameterwerte berücksichtigt wird.
\end{conceptbox}

\section{MAP und Regularisierung}

MAP ist wichtig, weil es direkt mit Regularisierung (regularization) zusammenhängt.

Wenn wir einen Prior verwenden, der große Parameterwerte unwahrscheinlich macht, dann werden solche Parameterwerte beim Schätzen bestraft.

Ein Beispiel ist ein Gauß-Prior auf einen Parametervektor \(\vec{w}\):

\[
    p(\vec{w})
    =
    \mathcal{N}
    \left(
        \vec{w}
        \mid
        \vec{0},
        \lambda^{-1}I
    \right).
\]

Dieser Prior bevorzugt kleine Gewichte.
Der Log-Prior enthält dann einen Term der Form:

\[
    -\lambda \vec{w}^{\top}\vec{w}.
\]

Beim Maximieren des Log-Posteriors entspricht das einer Bestrafung großer Gewichte.
Beim Minimieren eines negativen Log-Posteriors wird daraus ein Regularisierungsterm:

\[
    \lambda \vec{w}^{\top}\vec{w}.
\]

\begin{conceptbox}
MAP-Schätzung kann als Maximum-Likelihood-Schätzung mit zusätzlichem Prior-Wissen verstanden werden.

In vielen Modellen führt ein Gauß-Prior auf Parameter zu \(L_2\)-Regularisierung.
Bei linearer Regression ergibt das Ridge Regression.
\end{conceptbox}

\section{Bayesian Estimator}

MLE und MAP liefern jeweils einen einzelnen besten Parameterwert:

\[
    \hat{\theta}_{\mathrm{ML}}
\]

oder

\[
    \hat{\theta}_{\mathrm{MAP}}.
\]

Bayesian Estimation geht einen Schritt weiter:
Statt nur einen Punkt zu schätzen, betrachtet man die ganze Posterior-Verteilung.

\begin{definitionbox}
Ein Bayesian Estimator verwendet die Posterior-Verteilung

\[
    p(\theta \mid \vec{x})
    =
    \frac{
        p(\vec{x}\mid\theta)p(\theta)
    }{
        p(\vec{x})
    }.
\]

Die Unsicherheit über \(\theta\) bleibt also als Verteilung erhalten.
\end{definitionbox}

Man kann schreiben:

\[
    \hat{\Theta}_N
    \sim
    p(\theta \mid \vec{x}).
\]

Das bedeutet:
Die Schätzung ist nicht nur ein Punkt, sondern wird durch eine Verteilung über mögliche Parameterwerte beschrieben.

\begin{conceptbox}
Bayesian Estimation modelliert Unsicherheit über Parameter explizit.

Das ist besonders nützlich, wenn nur wenige Daten vorhanden sind oder wenn Unsicherheit in Entscheidungen einfließen soll.
\end{conceptbox}

\section{Vergleich: MLE, MAP und Bayesian}

\[
\begin{array}{c|c|c}
\text{Methode}
&
\text{Was wird berechnet?}
&
\text{Verwendet Prior?}
\\
\hline
\text{MLE}
&
\hat{\theta}_{\mathrm{ML}}
=
\arg\max_{\theta}p(\vec{x}\mid\theta)
&
\text{Nein}
\\
\hline
\text{MAP}
&
\hat{\theta}_{\mathrm{MAP}}
=
\arg\max_{\theta}p(\theta\mid\vec{x})
&
\text{Ja}
\\
\hline
\text{Bayesian}
&
p(\theta\mid\vec{x})
&
\text{Ja}
\end{array}
\]

\begin{notebox}
MLE und MAP geben Punktschätzungen.

Bayesian Estimation gibt eine ganze Verteilung über Parameter.
\end{notebox}

\section{Effizienz eines Schätzers}

Ein Schätzer ist effizient (efficient), wenn er im Vergleich zu anderen Schätzern einen kleinen erwarteten Fehler hat.
Meistens betrachtet man dafür den MSE.

\begin{definitionbox}
Ein Schätzer \(\hat{\theta}^{A}_N\) ist effizienter als ein Schätzer \(\hat{\theta}^{B}_N\), wenn gilt:

\[
    \mathrm{MSE}[\hat{\theta}^{A}_N]
    <
    \mathrm{MSE}[\hat{\theta}^{B}_N].
\]
\end{definitionbox}

\begin{notebox}
Effizienz hängt vom gewählten Fehlermaß ab.
In diesem Kurs ist der MSE das zentrale Fehlermaß für Schätzer.
\end{notebox}

\section{Zusammenhang zu maschinellem Lernen}

Schätzer sind eine Grundlage für viele ML-Algorithmen.

\begin{itemize}
    \item In der linearen Regression schätzen wir Gewichte \(\vec{w}\).
    \item In der logistischen Regression schätzen wir Parameter eines Klassifikators.
    \item In Gaussian Mixture Models schätzen wir Mittelwerte, Kovarianzen und Mischgewichte.
    \item Im EM-Algorithmus wechseln sich Posterior-Berechnung und Parameterschätzung ab.
    \item In neuronalen Netzen schätzen wir sehr viele Gewichte durch Optimierung einer Verlustfunktion.
\end{itemize}

\begin{conceptbox}
Viele Lernalgorithmen können als Parameterschätzung verstanden werden:

\[
    \text{Daten}
    \longrightarrow
    \text{Schätzer}
    \longrightarrow
    \text{Parameter}
    \longrightarrow
    \text{Modell}.
\]
\end{conceptbox}

\section{Häufige Fehler}

\begin{examtipbox}
Typische Fehler bei Schätzern:

\begin{enumerate}
    \item Parameter \(\theta\), Schätzer \(\hat{\theta}_N\) und Schätzwert \(\hat{\theta}_N(\vec{x})\) verwechseln.
    \item Vergessen, dass ein Schätzer eine Zufallsvariable ist.
    \item Unbiased mit gut gleichsetzen.
    \item Bias und Varianz getrennt betrachten, aber den MSE vergessen.
    \item Bei der Stichprobenvarianz durch \(N\) teilen, obwohl ein unverzerrter Schätzer gefragt ist.
    \item MLE für automatisch unverzerrt halten.
    \item Likelihood \(p(\vec{x}\mid\theta)\) mit Posterior \(p(\theta\mid\vec{x})\) verwechseln.
    \item Bei MAP den Prior vergessen.
    \item Bei der Log-Likelihood das Produkt nicht korrekt in eine Summe umwandeln.
\end{enumerate}
\end{examtipbox}

\section{Mini-Aufgaben}

\subsection{Aufgabe 1: Bias eines einfachen Schätzers}

Sei \(\hat{\theta}_N\) ein Schätzer für \(\theta\).
Angenommen:

\[
    \mathbb{E}[\hat{\theta}_N]
    =
    \theta + 2.
\]

Berechne den Bias.

\begin{examplebox}
Nach Definition:

\[
    \mathrm{Bias}[\hat{\theta}_N]
    =
    \mathbb{E}[\hat{\theta}_N]
    -
    \theta.
\]

Einsetzen ergibt:

\[
    \mathrm{Bias}[\hat{\theta}_N]
    =
    (\theta+2)-\theta
    =
    2.
\]
\end{examplebox}

\subsection{Aufgabe 2: MSE aus Bias und Varianz}

Angenommen:

\[
    \mathrm{Bias}[\hat{\theta}_N]
    =
    -3,
\]

und

\[
    \mathrm{Var}[\hat{\theta}_N]
    =
    5.
\]

Berechne den MSE.

\begin{examplebox}
Mit der Bias-Varianz-Zerlegung:

\[
    \mathrm{MSE}[\hat{\theta}_N]
    =
    \mathrm{Bias}[\hat{\theta}_N]^2
    +
    \mathrm{Var}[\hat{\theta}_N].
\]

Also:

\[
    \mathrm{MSE}[\hat{\theta}_N]
    =
    (-3)^2 + 5
    =
    9+5
    =
    14.
\]
\end{examplebox}

\subsection{Aufgabe 3: Stichprobenmittel}

Gegeben sind die Daten:

\[
    x_1=2,
    \quad
    x_2=4,
    \quad
    x_3=6.
\]

Berechne das Stichprobenmittel.

\begin{examplebox}
Das Stichprobenmittel ist:

\[
    \hat{\mu}_3
    =
    \frac{1}{3}
    (2+4+6)
    =
    4.
\]
\end{examplebox}

\subsection{Aufgabe 4: Unverzerrte Stichprobenvarianz}

Für dieselben Daten

\[
    x_1=2,
    \quad
    x_2=4,
    \quad
    x_3=6
\]

ist

\[
    \hat{\mu}_3 = 4.
\]

Berechne die unverzerrte Stichprobenvarianz.

\begin{examplebox}
Die unverzerrte Stichprobenvarianz ist:

\[
    \hat{s}^2_3
    =
    \frac{1}{3-1}
    \sum_{n=1}^{3}
    (x_n-\hat{\mu}_3)^2.
\]

Einsetzen:

\[
    \hat{s}^2_3
    =
    \frac{1}{2}
    \left[
        (2-4)^2
        +
        (4-4)^2
        +
        (6-4)^2
    \right].
\]

Also:

\[
    \hat{s}^2_3
    =
    \frac{1}{2}
    (4+0+4)
    =
    4.
\]
\end{examplebox}

\subsection{Aufgabe 5: MLE für Bernoulli}

Gegeben sind binäre Daten:

\[
    x_1=1,
    \quad
    x_2=0,
    \quad
    x_3=1,
    \quad
    x_4=1.
\]

Berechne den Maximum-Likelihood-Schätzer für \(\pi\).

\begin{examplebox}
Für Bernoulli-Daten gilt:

\[
    \hat{\pi}_{\mathrm{ML}}
    =
    \frac{1}{N}
    \sum_{n=1}^{N}
    x_n.
\]

Hier ist \(N=4\) und

\[
    \sum_{n=1}^{4}x_n
    =
    1+0+1+1
    =
    3.
\]

Also:

\[
    \hat{\pi}_{\mathrm{ML}}
    =
    \frac{3}{4}
    =
    0{,}75.
\]
\end{examplebox}

\section{Zusammenfassung}

\begin{itemize}
    \item Ein Parameter \(\theta\) ist eine feste, unbekannte Modellgröße.
    \item Ein Schätzer \(\hat{\theta}_N\) ist eine Funktion der Daten.
    \item Ein Schätzwert ist der konkrete Wert des Schätzers auf einem konkreten Datensatz.
    \item i.i.d. bedeutet unabhängig und identisch verteilt.
    \item Bias misst systematische Abweichung:
    \[
        \mathrm{Bias}[\hat{\theta}_N]
        =
        \mathbb{E}[\hat{\theta}_N]-\theta.
    \]
    \item Varianz misst die Schwankung des Schätzers:
    \[
        \mathrm{Var}[\hat{\theta}_N]
        =
        \mathbb{E}\left[
            \left(
                \hat{\theta}_N
                -
                \mathbb{E}[\hat{\theta}_N]
            \right)^2
        \right].
    \]
    \item Der mittlere quadratische Fehler ist:
    \[
        \mathrm{MSE}[\hat{\theta}_N]
        =
        \mathbb{E}\left[
            (\hat{\theta}_N-\theta)^2
        \right].
    \]
    \item Es gilt:
    \[
        \mathrm{MSE}
        =
        \mathrm{Bias}^2
        +
        \mathrm{Varianz}.
    \]
    \item Das Stichprobenmittel ist unverzerrt:
    \[
        \mathrm{Bias}[\hat{\mu}_N]=0.
    \]
    \item Die Varianz des Stichprobenmittels ist:
    \[
        \mathrm{Var}[\hat{\mu}_N]
        =
        \frac{\sigma_X^2}{N}.
    \]
    \item Die naive Stichprobenvarianz ist verzerrt:
    \[
        \mathrm{Bias}[\hat{\sigma}^2_N]
        =
        -\frac{1}{N}\sigma_X^2.
    \]
    \item Die unverzerrte Stichprobenvarianz ist:
    \[
        \hat{s}^2_N
        =
        \frac{1}{N-1}
        \sum_{n=1}^{N}
        (X_n-\hat{\mu}_N)^2.
    \]
    \item Maximum-Likelihood schätzt Parameter durch Maximieren der Datenwahrscheinlichkeit:
    \[
        \hat{\theta}_{\mathrm{ML}}
        =
        \arg\max_{\theta}
        p(\vec{x}\mid\theta).
    \]
    \item MAP maximiert den Posterior:
    \[
        \hat{\theta}_{\mathrm{MAP}}
        =
        \arg\max_{\theta}
        p(\theta\mid\vec{x}).
    \]
    \item Bayesian Estimation betrachtet die ganze Posterior-Verteilung \(p(\theta\mid\vec{x})\).
\end{itemize}

\section{Typische Prüfungsfragen}

\begin{examtipbox}
Mögliche Prüfungsfragen zu diesem Kapitel:

\begin{enumerate}
    \item Was ist der Unterschied zwischen Parameter, Schätzer und Schätzwert?
    \item Warum ist ein Schätzer eine Zufallsvariable?
    \item Was bedeutet i.i.d.?
    \item Definiere Bias eines Schätzers.
    \item Definiere Varianz eines Schätzers.
    \item Definiere MSE und leite die Bias-Varianz-Zerlegung her.
    \item Zeige, dass das Stichprobenmittel erwartungstreu ist.
    \item Berechne die Varianz des Stichprobenmittels.
    \item Warum ist die naive Stichprobenvarianz verzerrt?
    \item Warum verwendet man bei der unverzerrten Stichprobenvarianz \(N-1\)?
    \item Was ist der Bias der naiven Stichprobenvarianz?
    \item Was ist Maximum-Likelihood-Schätzung?
    \item Warum verwendet man die Log-Likelihood?
    \item Leite den MLE für eine Bernoulli-Verteilung her.
    \item Was ist der Unterschied zwischen MLE und MAP?
    \item Welche Rolle spielt der Prior bei MAP?
    \item Warum ist Bayesian Estimation mehr als eine Punktschätzung?
    \item Erkläre den Bias-Varianz-Trade-off.
    \item Warum ist ein unverzerrter Schätzer nicht automatisch der beste Schätzer?
    \item Wie hängt Parameterschätzung mit maschinellem Lernen zusammen?
\end{enumerate}
\end{examtipbox}